\clearpage\chapter{Basic Trigonometric Functions}
Like any other function, trigonometric functions can take inputs, produce outputs, be used to create a graph, transform/translate, etc. However, trigonometric functions are different in that they require additional vocabulary and properties to discuss.

The basic functions you'll learn in this chapter are the sine, cosine, and tangent functions. These three are the SOH-CAH-TOA you may have learned back in geometry. You'll be using them a \textit{lot} more here in pre-calculus and in all future calculus courses.

Let's start with sine.
\clearpage
\section{The Sine Function}
The sine function is often defined by the ratio
\[
\frac{opposite}{hypotenuse}
\]
where the opposite and hypotenuse are sides of a right triangle. It's also often written like this:
\[
\sin(\theta) = \frac{o}{h} \text{.}
\]
While this \textit{is} is a good explanation of it, I'll be diving deeper into what exactly this function is and how to graph it.